\documentclass[11pt,a4paper]{article}
\usepackage[T1]{fontenc}
\usepackage[utf8]{inputenc}
\usepackage{amsmath,amssymb,xcolor}
\usepackage[margin=25mm]{geometry}
\usepackage[colorlinks=true,urlcolor=blue]{hyperref}
\definecolor{nexorablue}{HTML}{2459B5}
\setlength{\parindent}{0pt}
\setlength{\parskip}{7pt}
\setlength{\emergencystretch}{2em}
\renewcommand{\familydefault}{\sfdefault}
\begin{document}

{\small\bfseries\color{nexorablue} NEXORA CPGE}\par
{\small MPSI / NEX-LOG-04 / Corrigé}\par
{\LARGE\bfseries Un maximum obtenu en un seul point}\par
\vspace{0.25em}\hrule\vspace{0.6em}

\section*{Existence et unicité}
On a $3x-x^2=9/4-(x-3/2)^2$. Atteindre au moins $9/4$ force le carré à valoir zéro. Le réel $3/2$ convient et il est unique.

\section*{Second seuil}
Le carré doit être au plus $1/4$, soit $|x-3/2|\le1/2$. On obtient le segment $[1,2]$, bornes incluses.


\end{document}
